% New proved local comparison after checkpoint09.  Root assigns the result ID.
% Its ambient branch-output realization belongs to the separate finite compiler.
\subsection{A collar correction retaining the diagonal four-point disk}

Let $T$ carry the actual involution $\iota:z\mapsto-z$, the marked
origin $O$, and four named pairs $\{p_j,-p_j\}$ of nonfixed points.
Write $S$ for this nine-point torus and $\overline S=T/\langle\iota\rangle$
for its actual branch sphere.  Its distinguished branch mark is
$b_0=\overline O$, its other three branch points form a set $B$, and
the four named pair images are $t_1,\ldots,t_4$.

Choose an embedded closed disk $D\subset\overline S$ whose interior
contains the four $t_j$ and whose closure misses all four branch
points.  One constructs it as a small regular neighbourhood of an
embedded tree joining the four $t_j$ in the branch-point complement.
The actual inverse image of $D$ is the disjoint union of two disks
$D_+,D_-$, exchanged by $\iota$; each restriction to $D$ is an
orientation-preserving homeomorphism.  Transport the same ordered
braid marking through these two maps.  Extension by the identity
outside their interiors gives the homomorphism
\[
 \begin{split}
 \Delta:B_4&\longrightarrow C_{\Gamma_{1,[9]}}(\iota),\\
 \sigma_i&\longmapsto\sigma_{i,+}\sigma_{i,-}
                                      \qquad(1\leq i\leq3).
 \end{split}                                                     \tag{DDP.1}
\]
The two factors commute because their supports are disjoint.  The
involution exchanges the factors, proving the displayed centralizer
membership under the actual surface map.

The ordinary branch quotient followed by forgetting $B$ gives, on
the subgroup fixing $O$ and each of the four named pairs,
\[
 \delta:C(\iota)\longrightarrow\Gamma_{0,5}.
                                                               \tag{DDP.2}
\]
Here $C(\iota)$ denotes that specified subgroup of the ordinary
centralizer.  The exact ordinary quotient map is the one in
(IFM.1)--(IFM.2): downstairs it retains $b_0,t_1,\ldots,t_4$ and
forgets the unordered set $B$.  In particular
\[
 \delta\Delta|_{P_4}=\pi:P_4\twoheadrightarrow
        P_4/\langle(\sigma_1\sigma_2\sigma_3)^4\rangle
                    =\Gamma_{0,5}.                             \tag{DDP.3}
\]
To check the map, descend a braid supported on $D_+\cup D_-$.
Its quotient is the original braid on $D$.  After the three branch
points are forgotten, the complementary disk has just the marked
point $b_0$.  Capping the outer braid boundary by that once-marked
disk kills exactly its full boundary twist.  This is the stated
central quotient, with every named $t_j$ retained.

\paragraph{The finite local comparison.}
For the finite action of A.4.40, use its actual groups $U_S(m)$ and
the corresponding finite five-point input $\mathcal P_1(m)$.
There is a starting quilt $Q$ and, for every actual
$F=(\lambda,f)\in\mathcal P_1(m)$, an explicitly specified pure
multitwist $c_F$ in the image of the complement of
$D_+\cup D_-$ such that
\[
 \Psi_F(\iota)=c_F\,q(\iota)c_F^{-1},\qquad
 [c_F,q\Delta(B_4)]=1.                                         \tag{DDP.4}
\]
Here $q$ is the actual finite quotient map.  In particular
$\Psi_F^0=\operatorname{Ad}(c_F^{-1})\Psi_F$ fixes the named
involution and retains exactly the same two disk-generator images.
With the three ordinary braid words
\[
 \begin{aligned}
 b_1(F)&=\sigma_1^\lambda,\\
 b_2(F)&=f(\sigma_1^2,\sigma_2^2)^{-1}
                  \sigma_2^\lambda f(\sigma_1^2,\sigma_2^2),\\
 y_3&=\sigma_2\sigma_1^2\sigma_2,\\
 b_3(F)&=f(y_3,\sigma_3^2)^{-1}
                  \sigma_3^\lambda f(y_3,\sigma_3^2),
 \end{aligned}                                                   \tag{DDP.5}
\]
the actual finite identities are
\[
          \Psi_F^0(q\Delta(\sigma_i))
                      =q\widehat\Delta(b_i(F))
                              \qquad(1\leq i\leq3).             \tag{DDP.6}
\]
The hats in this display mean evaluation of the displayed profinite
words through the specified ordinary braid homomorphism and then
through the finite map $q$.  They do not assert an unproved infinite
$B_4$ action for every finite-pentagon parameter.

\paragraph{Proof of the collar correction.}
The complement $C=T\setminus(\operatorname{int}D_+\cup
\operatorname{int}D_-)$ is a genus-one surface with two exchanged
boundary components and the marked origin.  Its quotient is a disk
with four branch points.  Inside that quotient disk choose a circle
surrounding $b_0$ and precisely one other branch point.  Its two
inverse-image circles $a_1,a_2$ are exchanged by $\iota$.
Their inside is an annulus containing $O$, hence a pair of pants
when $O$ is blown up.  Call that pant $P_0$.
In the complementary annulus downstairs, join the remaining two
branch points by an embedded arc.  Its inverse image is a circle
$a_3$.  Cutting there gives two pants $P_+,P_-$ exchanged by the
involution.  Their external cuffs are $\partial D_+$ and
$\partial D_-$.  Thus the exact core pants graph has the fixed
pant $P_0$, the exchanged pants $P_+,P_-$, exchanged cuffs
$a_1,a_2$, and the fixed cuff $a_3$.

Blow up all nine marked points to parametrized boundary circles.
Choose the eight paired parametrizations as actual involution
translates.  Near $O$, the involution is a half-turn.  Changing it
in the origin collar to fix the new boundary pointwise gives a
literal lift $\widetilde\iota$ with
\[
                  \widetilde\iota^2=t_0^{\epsilon},
                       \qquad\epsilon\in\{1,-1\},              \tag{DDP.7}
\]
where $t_0$ is the positive origin-boundary twist.  The sign is
determined by the chosen collar interpolation.  Indeed the square
is the full turn of that collar and is identity off the collar.
Replacing the lift by $t_0^{-\epsilon}\widetilde\iota$ changes the
sign; neither choice changes its marked cap to $\iota$.

Take a pants decomposition consisting of the three core pants and
the two transported identical disk decompositions.  Choose the
quilt on $D_-$ to be the actual translate of the quilt on $D_+$,
and choose the two core quilts on $P_+,P_-$ likewise.  Match the
remaining seam endpoints through the core collars.  A pant quilt
here is exactly the seam cellulation of \cite[\S7]{ET_NS2000}: one
seam joins each pair of its three actual cuffs, and the seam set
cuts the pant into two patches.  The quilt object has no additional
ordering of those three seams or of the two patches.  Therefore the
transposition of the two non-origin cuffs of $P_0$ sends its quilt
to a quilt over the same pant, with the cuff labels transported by
the actual homeomorphism.  It introduces no extra permutation or
R-edge.  Cutting the two seam cellulations into hexagons gives
exactly the collar displacement coordinates already used in
(TA.9) and (FGA.7).

The involution preserves this pants decomposition.  Thus its quilt displacement
uses only the actual half-adjustments of (TA.9) and (FGA.7)--(FGA.9),
with a unique integer vector $N=(N_a)$:
\[
        \widetilde\iota^{-1}Q
               =\left(\prod_a D_a^{N_a/2}\right)Q.              \tag{DDP.8}
\]
The coordinates on cuffs internal to $D_+$ or $D_-$ are zero:
the two restricted quilts are exact involution translates, and the
free collar-coordinate description restricts to those same disk
seam cellulations.  Possible adjustments at their outer cuffs are
retained as core collar adjustments.

Let $P$ permute collar coordinates by the involution.  Orientation
preservation sends each positive twist to the positive twist on
the image cuff.  Apply $\widetilde\iota^{-1}$ twice to (DDP.8)
and use (DDP.7).  Freeness of the integer half-adjustment lattice
gives the exact equation
\[
                       N+PN=-2\epsilon e_0.                    \tag{DDP.9}
\]
In particular every fixed-cuff coordinate other than the origin
is zero, and the two coordinates on a nontrivial orbit are
opposites.  This step is an equality of integer seam displacements,
before any finite quotient is taken; it is not obtained by dividing
a finite even-order congruence.

Put $\nu=(\lambda-1)/2$.  Choose one cuff $a$ from each two-element
orbit of the core collar set, and define
\[
                     c_F=\prod_{\{a,\iota a\}}q(t_a)^{N_a\nu}.
                                                               \tag{DDP.10}
\]
All factors commute.  The half-adjustment path in (DDP.8) has weight
$\prod_aq(t_a)^{N_a\nu}$.  Marked capping kills $t_0$ and retains
the other collar maps, by A.4.40.  After that cap the weight is
exactly $c_F\,{}^{q(\iota)}c_F^{-1}$.  Substituting it into
(FGA.9) proves the first identity in (DDP.4).  Every factor in
(DDP.10) is supported in the core or an outer disk collar, so it
commutes with the two disk braid groups.  This proves the second
identity.  In a finite target each power in (DDP.10) is the value
of an ordinary integer power; consequently $c_F$ is in the actual
geometric pure image of that core.  The parameter $\nu$ only needs
to be retained modulo the exponent of the finite target, so the
same factor-two multiplier precision as (FPA.13)--(FPA.14) suffices.
Taking the least nonnegative representative of that residue of
$\nu$ in (DDP.10) supplies one literal ordinary multitwist word.

There is no A- or S-move in (DDP.8).  Thus this involution correction
introduces no extra $f$ evaluation and no extra coefficient test.
Choose on $D_+$ the standard genus-zero braid quilt; transport it
to $D_-$ as above.  The ordinary sewing maps of A.4.40 give the
same words (DDP.5) on both disk factors.  Their disjoint supports
make their product precisely the image under $\Delta$ of that
same word.  Since the correction commutes with both factors, this
proves (DDP.6).

\paragraph{The exact parameter record on the final five-point generators.}
Every pure generator $A_{ij}$ of $P_4$ is a specified word in the
three $\sigma_i$.  Substitute (DDP.5) in that word and call the
result $A_{ij}(F)$.  Each conjugating $f$ word in (DDP.5) is pure,
and $\lambda$ is odd, so the braid permutation of each $\sigma_i$
is unchanged.  Thus $A_{ij}(F)$ lies in the completed pure subgroup.
The marked homomorphisms (DDP.1)--(DDP.3)
give, by evaluation through their profinite completions,
\[
 \widehat\delta\bigl(\widehat\Delta(A_{ij}(F))\bigr)
                         =\widehat\pi(A_{ij}(F)).                \tag{DDP.11}
\]
The right side is the standard same-parameter five-point generator
word.  Equations (DDP.6) and (DDP.11) provide the generatorwise
comparison for the retained finite branch-and-forgetting diagram.
They use the actual surjection (DDP.3); no assertion that four
other five-point quotients determine an eight-point automorphism
is used.  Realizing the specified branch-output quotient inside an
ambient finite mapping-class image remains the separate ambient
construction of that diagram; the local proof here does not extend
an arbitrary subgroup quotient without that construction.

This construction uses the nine-point torus, with only $O$ marked
among the four involution fixed points.  The ordinary isomorphism
(IFM.1) transports its centralizer maps to the surface with the other
three fixed points retained as an unordered block.  No invariant
pants decomposition after those three additional markings is asserted
or needed for (DDP.1)--(DDP.11).
